\documentclass[11pt,a4paper]{article}
\usepackage{fontspec}
\setmainfont{Noto Serif}
\setsansfont{Noto Sans}
\usepackage{babel}
\babelprovide[main,import]{russian}
\usepackage{amsmath,amssymb,mathtools,booktabs,geometry,fancyhdr,microtype,xcolor,hyperref}
\geometry{left=21mm,right=21mm,top=21mm,bottom=21mm}
\definecolor{ink}{HTML}{173D57}
\hypersetup{colorlinks=true,urlcolor=ink,linkcolor=ink,pdftitle={BGS-F. Возвраты со сдвигом}}
\pagestyle{fancy}\fancyhf{}\lhead{\small BGS-F}\rhead{\small Возвраты со сдвигом}\cfoot{\thepage}
\setlength{\headheight}{15pt}\setlength{\parindent}{1em}\setlength{\parskip}{2.5pt}
\setlength{\emergencystretch}{3em}\tolerance=3000
\newcommand{\Tr}{\operatorname{Tr}}
\newcommand{\EE}{\mathbb E}
\newcommand{\R}{\mathbb R}
\newcommand{\Z}{\mathbb Z}
\newcommand{\C}{\mathbb C}
\newcommand{\si}{\mathrm{Sine}_2}
\newcommand{\eps}{\varepsilon}
\newcommand{\sinc}{\operatorname{sinc}}
\newcommand{\sect}[1]{\section*{\sffamily\color{ink}#1}}
\newcommand{\subt}[1]{\par\medskip\noindent\textbf{#1.}\quad}
\newcommand{\T}{\mathbb T}
\newcommand{\dist}{\operatorname{dist}}
\newcommand{\Ad}{\operatorname{Ad}}
\begin{document}
\noindent{\small\sffamily\color{ink} МАТЕМАТИЧЕСКАЯ ФИЗИКА\hfill 20 НОЯБРЯ 2026}
\begin{center}{\LARGE\sffamily Подавление возвратов со сдвигом}\\[4pt]
{\large Исходный кик и равномерные оценки по всем степеням}
\end{center}
Пусть $U_N=\exp(i10^{-80}NA_N)M_N$ --- исходный оператор, $T_N(v)$ --- его модульные операторы Вейля, а $W_N=\operatorname{Ad}U_N$ --- сопряжение на пространстве матриц с нормой $\|X\|_{\rm n}^2=N^{-1}\Tr X^*X$. Для центрированных представителей $v=(v_1,v_2)$ положим $n(v)=|v_1|+\sqrt2|v_2|$. При любой фиксированной $\delta>0$ доказано
\[
 \boxed{\displaystyle
  \sup_{m\in\Z}\ \sup_{n(v)\ge\delta N}
      \frac{|\Tr(T_N(v)^*U_N^m)|}{N}
      =O_\delta\!\left((\log\log\log N)^{-1/2}\right).}                \tag{1}
\]
Оценка справедлива по всем достаточно большим нечётным $N$. В частности, ни одна степень $U_N$ не приближается в нормированной норме Гильберта--Шмидта к скалярному кратному макроскопического вейлева сдвига. Минимальное расстояние по скалярным фазам равномерно стремится к $\sqrt2$.

Основная оценка получена для растущего числа итераций сопряжения:
\[
 \sup_{n(v)\ge\delta N}\ \sup_w |(W_N^j)_{wv}|
     \le e^{-(\log N)^{3/4}},\qquad
               1\le j\le T=\lfloor c\log\log\log N\rfloor.           \tag{2}
\]
Здесь $c>0$ можно фиксировать, а порог по $N$ зависит от $\delta$. Ранее степенная оценка была получена при каждом фиксированном $j$. Для (2) понадобились равномерная геометрическая отделённость всех классических итераций от перевода сдвига в сдвиг, оценка объёма малых отклонений и контроль констант квантового переноса.

Качественный вывод охватывает все $v\ne0$, включая сдвиги размера $1/N$. При $N=7^k$ и $m=3N/49$ доказано также $\min_{|\zeta|=1}\|U_N^m-\zeta M_N^m\|_{\rm n}\ge\sqrt{12/7}-o(1)$. У линейной степени в этом случае 49 вейлевых коэффициентов, из которых 48 подавляются одновременно. Это даёт разрыв от частичного линейного возврата при времени порядка $N$.

Из (2) выводится аналог (1) для целых пакетов из $\exp(\tfrac12(\log N)^{3/4})$ внешних мод и для любого оператора, коммутирующего с $U_N$. На размерностях $N=7^j$, $j\ge2$, это даёт асимптотическое равенство 49 торсионных следов $\Tr(Q_\chi U_N^m)$ с ошибкой $O(N/\sqrt T)$, равномерно по $m$. Получено также равенство соответствующих весов в локальном точечном процессе при равномерном выборе начала развёртки.

\textbf{Статус.} Оценки (1),(12) исключают возвраты к ненулевому сдвигу. Нулевая мода даёт обычный след $\Tr U_N^m$. Его вклад на гейзенберговской полосе, включая скалярные арифметические времена, остаётся открытым.
\newpage
\sect{1. Отделённость от сдвигов при всех итерациях}
Физическая карта есть $F=\phi^\tau\circ B$, $B=\left(\begin{smallmatrix}3&-4\\-2&3\end{smallmatrix}\right)$, где $\phi$ --- гамильтонов поток для физического символа (2.5) исходной постановки. Пишем $H^{-1}=\mathrm{id}+k$ для сопряжения $FH=HB$ на торе. Из части о жёсткости Вейля, разделы 5 и 10, используем следующие точные утверждения для фиксированного кика:
\[
 \|H-\mathrm{id}\|_\infty,\ \|k\|_\infty\le100\tau,\quad
 |\widehat{\ell_u k}(B^{\mathsf T}e_i)|>10^{-4}\tau\quad (i=1,2),
 \quad F^j t_a F^{-j}=t_b\ \Longrightarrow\ a=b=0.                   \tag{3}
\]
Здесь $\ell_u=(1/2,-1/\sqrt2)$, $\ell_u B=(3+2\sqrt2)\ell_u$, а $t_a$ есть трансляция на $a$. Обе выделенные частоты образуют целочисленный базис.

Для $\Phi_a=H^{-1}t_aH$ поднятие её смещения при $x=H^{-1}(y)$ равно $a+k(y+a)-k(y)$. Если $\operatorname{dist}(a,0)\ge\delta_0>0$, хотя бы один из двух множителей $|e^{2\pi i(B^{\mathsf T}e_i)\cdot a}-1|$ ограничен снизу положительной константой, равномерно по $a$. Поэтому осцилляция $\ell_u(\Phi_a(x)-x)$ не меньше $d_{\delta_0}\tau>0$ по (3). Осцилляция той же компоненты у поднятия $B^j\Phi_aB^{-j}-\mathrm{id}$ не меньше $(3+2\sqrt2)^j d_{\delta_0}\tau$.

Предположим, что $F^j t_a F^{-j}$ в равномерной метрике тора находится на расстоянии меньше $1/100$ от $t_b$. Тогда $B^j\Phi_a B^{-j}$ находится на расстоянии не больше $1/100+200\tau$ от той же трансляции. Достаточно применить $H$ ко входу и $H^{-1}$ к выходу, каждый из них отклоняет точку не более чем на $100\tau$. В поднятии разность отличается от фиксированного целого вектора на величину не более $1/100+200\tau$: ближайший целый вектор единственен и непрерывен по входу, следовательно постоянен. Осцилляция его неустойчивой компоненты тогда меньше 1. Это невозможно при $j\ge j_0(\delta_0,\tau)$. Оставшиеся конечные $j$ контролируются по компактности параметров $a,b$ и отсутствию точного равенства из (3).

Итак, существует $\eta_{\delta_0}>0$ такое, что одновременно для всех $j\ge1$, всех $a$ с $\operatorname{dist}(a,0)\ge\delta_0$ и всех $b$
\[
   \sup_{x\in\T^2}\operatorname{dist}
             (F^j(x+a)-F^j(x)-b,0)\ge\eta_{\delta_0}.                 \tag{4}
\]
Здесь и далее расстояние на торе берётся по максимальной норме. Уже в (4) число итераций может расти. Тот же вывод для $G=t_{-b}F^jt_a F^{-j}$ получается заменой переменной $x$ на $F^jx$.
\newpage
\sect{2. Объём области малых смещений}
Обоснуем количественный переход от (4) к оценке объёма. Поднятие $F$ вещественно-аналитично, имеет вид $Bx$ плюс периодическая часть и продолжимо на трубку $|\operatorname{Im}z|<r_0$. На трубке меньшего фиксированного радиуса $\|DF\|\le K$, $K>1$, так что $|\operatorname{Im}F(z)|\le K|\operatorname{Im}z|$. Значит, $F^j$ продолжимо на трубку радиуса $r_j=r_0K^{-j}$ и мнимая часть выхода там ограничена $r_0$. Координатные функции
\[
 f_l(z)=\exp\{2\pi i(F^j(z+a)-F^j(z)-b)_l\}-1,\qquad l=1,2,
\]
имеют на этой трубке равномерную константу $M$ для модуля. Для каждого $j,a,b$ по (4) есть $l$ и вещественная точка $x_0$ с $|f_l(x_0)|\ge\eta'>0$, где $\eta'$ единообразна по параметрам из (4).

Рассмотрим произвольный вещественный единичный вектор $u$ и $f_l(x_0+zu)$ как ограниченную голоморфную функцию в полосе $|\operatorname{Im}z|<r_j$. Для $0\le z\le2$ преобразование $\zeta=\tanh(\pi z/(4r_j))$ переводит полосу в единичный диск и отрезок в $[0,R]$, $R=\tanh(\pi/(2r_j))$. Положим $s=1-R$, $R_1=R+s/3$, $R_2=R+2s/3$. Число нулей в диске радиуса $R_1$ не превосходит $D=C/s$ по формуле Йенсена на радиусе между $R_2$ и 1, поскольку значение в нуле не меньше $\eta'$. Постоянная $C$ зависит от $M/\eta'$. Радиус $R_1$ можно немного изменить, сохранив сравнимые зазоры и избегая нулей на границе. Удалим эти нули их конечным произведением Бляшке. Остаток $h$ ограничен $M$ в диске (по максимальному принципу, переходя к радиусам, стремящимся к 1), без нулей до $R_1$, и $|h(0)|\ge\eta'$. Неравенство Харнака для $\log(M/|h|)$ даёт $|h|\ge M e^{-C'/s}$ на $[0,R]$.

Каждый фактор Бляшке с нулём $\zeta_i$ по модулю не меньше $|\zeta-\zeta_i|/2$. Если исходный модуль не больше $\varepsilon$, то хотя бы один из нулей расположен на расстоянии $O(\varepsilon^{1/D})$ от $\zeta$ при $\varepsilon<M e^{-C'/s}$. Это следует из произведения факторов и $(M e^{-C'/s})^{-1/D}=O(1)$; для нулевого числа нулей множество пусто. В обратной координате производная не больше $C/s$. Итак, длина множества малых значений на вещественном отрезке не больше $C s^{-2}\varepsilon^{1/D}$; та же граница с увеличенной константой верна и при больших $\varepsilon$. Полярные координаты вокруг $x_0$ покрывают фундаментальную ячейку радиусом 2; радиальный множитель ограничен. Условие малого торического смещения влечёт $|f_l|\le 2\pi\varepsilon$. Поскольку $1/s\le C\exp(C K^j)$, получили
\[
 \operatorname{area}\{x:\operatorname{dist}(G(x),x)<\varepsilon\}
      \le C_\delta \exp(C_\delta K^j)
                   \varepsilon^{\exp(-C_\delta K^j)}.              \tag{5}
\]
Константы увеличены, чтобы учесть округление $D$ и фиксированную верхнюю границу радиуса. Источник формулы Йенсена и представления факторов приведён в конце.
\newpage
\sect{3. Рост констант квантового переноса}
Для $v,w$ взяты точные сеточные сдвиги $a=(-v_1,v_2)/N$, $b=(-w_1,w_2)/N$ на торе. При $n(v)\ge\delta N$ первый отделён от нуля. Коэффициент $(W_N^j)_{wv}$ есть нормированный след унитарного оператора
\[
                    V=T_N(w)^*U_N^j T_N(v)U_N^{-j}
\]
с классической картой $G=t_{-b}F^jt_a F^{-j}$. Усредним $T_N(k)^* V T_N(k)$ по $0\le k_1,k_2<L$. След не меняется, а квадрат его нормированного модуля ограничен квадратом нормированной гильбертовой нормы усреднения, то есть $\|\frac1{L^2}\sum_k T_N(k)^* V T_N(k)\|_{\rm n}^2$. В разложении этой нормы с суммарным весом один возникают корреляции $N^{-1}\Tr(V^* T_N(k-k') V T_N(k-k')^*)$. Фазы проектных операторов Вейля сокращаются.

Укажем константы остатка. Для одного кика и $f\in C^7(\T^2)$ ошибка Егорова в операторной норме не больше $C N^{-2}\|f\|_{C^7}$. Коммутатор в физическом базисе $W(a,b)=(-1)^{ab}T_N(a,b)$ отличается от квантования скобки на величину $\le C N^{-2}\sum_l(1+|l|)^3|\widehat f(l)|$ после деления на $\hbar$. Здесь используется оценка кубического остатка синуса и конечный носитель символа кика; ряд оценивается нормой $C^7$. Интегрирование по времени фиксированного потока даёт указанную ошибку для полного кика. Линейная часть и сеточные сдвиги действуют точно. Применение к обратному шагу даёт такую же оценку, так как гамильтонов поток обратим. Оператор переноса $f\mapsto f\circ F^{\pm1}$ ограничен в $C^7$ с фиксированной нормой. Поэтому телескопирование по $2j$ шагам даёт ошибку $\le C^{j+1} N^{-2}(1+L)^7$ для характера порядка не выше $2L$. Все производные сдвигов остаются единичными.

Произведение двух квантований отличается от квантования произведения не более чем на $C N^{-1}\|f\|_{C^5}\|g\|_{C^5}$ по точной фазе Вейля. Ошибка нормированного следа квантованного произведения относительно интеграла не больше $C N^{-4}\|fg\|_{C^4}$ по сумме частот, кратных $N$. В совокупности, с увеличением постоянной $C$, остаток в корреляциях не превосходит $C^{j+1}N^{-1}(1+L)^{12}$. В интеграле перенос обратной картой можно заменить переносом прямой, меняя переменную по сохранению площади.

Суммируя классические корреляции, получаем интеграл произведения двух квадратов нормированных геометрических сумм длины $L$, на компонентах $G(x)-x$. Вне окрестности $L^{-1/2}$ хотя бы один квадрат не больше $C/L$. Внутри применяем (5). При $L=\lfloor N^{1/24}\rfloor$ получаем
\[
 |(W_N^j)_{wv}|^2
 \le C_\delta e^{C_\delta K^j}
                L^{-\tfrac12e^{-C_\delta K^j}}+
                 C/L + C^{j+1} N^{-1}(1+L)^{12}.                   \tag{6}
\]
Выберем $c>0$ так, чтобы $c\log K<1/2$. Тогда для $j\le c\log\log\log N$ знаменатель показателя в первом члене не превосходит $(\log N)^{1/8}$ при достаточно большом $N$. Последние два члена убывают степенно по $N$. Первый оценивается через $\exp(-c'(\log N)^{7/8})$, что доказывает (2).
\newpage
\sect{4. Спектральные проекторы и произвольная степень}
Пусть $E$ состоит из внешних меток $n(v)\ge\delta N$ и
$|E|\le\exp(\tfrac12(\log N)^{3/4})$. Пусть $P_E$ --- ортогональный проектор на их линейную оболочку в пространстве матриц. Для $1\le j<T$ имеем
\[
        \|P_E W_N^j P_E\|
           \le |E| e^{-(\log N)^{3/4}}
           \le e^{-\tfrac12(\log N)^{3/4}}
\]
по оценке максимального элемента, умноженного на размер. Для единичного $X$ из этой оболочки рассмотрим среднее
\[
         A_\theta X=\frac1T\sum_{j=0}^{T-1} e^{-ij\theta}W_N^jX .
\]
Раскрытие квадрата нормы даёт $\|A_\theta X\|_{\rm n}^2\le T^{-1}+e^{-\tfrac12(\log N)^{3/4}}$, равномерно по $\theta$. По спектральной теореме множитель фильтра ограничен снизу положительной абсолютной константой в дуге $|\varphi-\theta|\le1/T$: все его слагаемые имеют вещественную часть не меньше $\cos 1$. Если $\Pi_J$ --- спектральный проектор $W_N$ на произвольную дугу $J$, покрытие такими малыми дугами даёт
\[
                         \|\Pi_J P_E\|^2\le C(|J|+T^{-1}).           \tag{7}
\]
Для одноточечного проектора $\Pi_{\{e^{i\theta}\}}$ получаем константу $1+o(1)$ при $T^{-1}$, так как среднее равно единице в центре фильтра. Наличие кратных собственных значений ни в (7), ни ниже не требует дополнительных условий.

При $\theta=0$ проектор $\Pi_{\{1\}}$ имеет образом коммутант $U_N$. Для любого $Z$ с $ZU_N=U_N Z$ по сопряжённости проекторов
\[
        \left(\sum_{v\in E}
               |N^{-1}\Tr(T_N(v)^*Z)|^2\right)^{1/2}
                          \le (1+o(1))T^{-1/2}\|Z\|_{\rm n}.        \tag{8}
\]
Берём $Z=U_N^m$ для любого целого $m$ или $Z=f_N(U_N)$ для произвольной ограниченной функции на его спектре. Это доказывает (1) и даёт оценку всего пакета. В частности, квадрат расстояния от $U_N^m$ до фазовых кратных $T_N(v)$ равен $2-2|\Tr(T_N(v)^*U_N^m)|/N$ при оптимальной фазе. Если $Y$ --- единичный в норме $\|\cdot\|_{\rm n}$ линейный пакет этих мод, то скалярное произведение $Y$ с $U_N^m$ также ограничено по (8). Таким образом, результат охватывает произвольные когерентные суперпозиции в заданном пакете.

Вывод по всем степеням основан на том, что $U_N^m$ является неподвижным вектором для сопряжения. В (2) оцениваются короткие, растущие времена \emph{сопряжения}; ограничение на $j$ не устанавливает никакого ограничения на $m$ в (1). У этой проекции всегда есть тождественная мода с единичным весом, поэтому обычный след этим методом не ограничен.
\newpage
\sect{5. Секторы на простых степенях и локальная развёртка}
Зафиксируем нечётное $d\ge3$ и возьмём $d^2\mid N$. Положим $h=N/d$. Операторы $T_N(hu)$ для $u\in(\Z/d\Z)^2$ дают представление конечной абелевой группы: по модульному правилу Вейля их фаза произведения равна 1, так как $N\mid h^2$. Проекторы на характеры
\[
             Q_\chi=\frac1{d^2}\sum_u \overline{\chi(u)} T_N(hu)
\]
ортогональны, дают единицу в сумме и имеют ранги $N/d^2$. Последнее следует из нулевых следов всех ненулевых сдвигов. При $d$ фиксированном все ненулевые метки отделены от нуля на макроскопическом масштабе, поэтому (1),(8) дают
\[
   \Tr(Q_\chi U_N^m)=d^{-2}\Tr U_N^m+O_d(N/\sqrt T)
\]
равномерно по целому $m$. Точнее, по Парсевалю с $p_m=\Tr U_N^m$ имеем
\[
 \begin{split}
    \frac1{N^2}\sum_\chi
          |\Tr(Q_\chi U_N^m)-p_m/d^2|^2
      &=\frac1{d^2}\sum_{u\ne0}
          |N^{-1}\Tr(T_N(hu)^*U_N^m)|^2 \\
      &\le\frac{1+o(1)}{d^2 T}.
 \end{split}                                                        \tag{9}
\]
В частности, при $N=7^j$, $j\ge2$, и $m=3rN/7$, $r=1,2,3,4$, знаковая сумма перекрёстных членов между конечными секторными следами равна
\[
       \sum_{\chi\ne\psi}\Tr(Q_\chi U_N^m)
                \overline{\Tr(Q_\psi U_N^m)}
                =\frac{48}{49}|p_m|^2+O(N^2/T).                     \tag{10}
\]
Ошибка в (10) неположительна и в точности равна минус сумме квадратов в (9) до нормировки. Здесь проектор ставится в след один раз; интерференция между полными историями секторных переходов может иметь другие вклады.

Для любого ортонормированного собственного базиса $(e_j)$ положим $w_\chi(j)=\langle e_j,Q_\chi e_j\rangle$. Диагональная проекция по этому базису лежит внутри коммутанта. Для каждого ненулевого $u$ её норма на $T_N(hu)$ не более $(1+o(1))/\sqrt T$ по (7),(8). Раскрывая $Q_\chi$, получаем
\[
           \frac1N\sum_j |w_\chi(j)-d^{-2}|^2=O_d(T^{-1}).            \tag{11}
\]
Это равномерно по выбору базиса также при кратностях. Для $\lambda_j=e^{2\pi i\theta_j}$ и равномерного $\alpha$ рассмотрим развёрнутый процесс с весами $w_\chi(j)$ на точках $N(\theta_j-\alpha+l)$, $l\in\Z$. Для любого теста $f\in L^1(\R)$ с компактным носителем среднее абсолютное отклонение его интеграла от $d^{-2}$ интеграла по невзвешенному процессу не больше
\[
       \frac{\|f\|_1}{N}\sum_j |w_\chi(j)-d^{-2}|
                             = O_d(\|f\|_1/\sqrt T).
\]
Интегрирование по $\alpha$ покрывает прямую в точности после суммы по $l$. Это утверждение относится к микроскопической развёртке и сохраняется при любом слабом локальном пределе исходного процесса.
\newpage
\sect{6. Что установлено для арифметических времён}
При $N=7^j$ формула (10) одинакова на всех четырёх временах скалярного возврата линейной карты в полной полосе $m<2N$. После деления на $N^2$ и взвешивания эти четыре ошибки имеют суммарный порядок $O(1/T)$. Таким образом, на масштабе возможного большого пика $|p_m|\asymp N$ его секторные следы не могут быть сосредоточены в собственном подмножестве из 49 секторов. Для исключения самого пика требуется контроль нулевой моды.

На пеллевых размерностях $N=h_k$, $k$ нечётно, использованная ранее оценка давления даёт $|\Tr U_N^k|/N=O(N^{-\eta})$ для некоторого $\eta>0$. В сочетании с (8) она подавляет любую проекцию $U_N^k$ на единицу и на пакет внешних мод допустимой мощности. Для сравнения, при нулевом кике $M_N^k$ скалярна. Это относится к первому пеллеву возврату; оценка нулевой моды на его растущих кратных ещё отсутствует. Для каждого фиксированного $d$ весь тот же вывод верен для секторных следов на пеллевых размерностях с $d^2\mid h_k$.

\textbf{Зависимости.} Классические коэффициенты сопряжения и отсутствие точного перехода между сдвигами взяты из разделов 5 и 10 части о жёсткости Вейля. Оценка для обычного следа на первом пеллевом возврате взята из части о давлении. Растущее окно (2), вывод (8) для всех степеней, формулы (9)--(11) и локальная развёртка выведены здесь.

\textbf{Конечные проверки.} В архиве приведён целочисленный контроль фаз формулы проектора для $d^2\mid N$, независимые проверки Парсеваля, равенства определённой части форм-фактора и оценки спектрального фильтра, а также диагностика весов на малых матрицах. Численные диагностики используют увеличенную силу кика для проверки тождеств. В доказательствах (1)--(11) сила в точности $10^{-80}$.

\textbf{Источник аналитических тождеств.} Формула Йенсена и конечные произведения Бляшке изложены в лекциях К. Макмаллена \emph{Complex Analysis}, разделы 44--45. \href{https://people.math.harvard.edu/~ctm/home/text/class/berkeley/205/93/course/course.pdf}{Версия автора}.
\newpage
\sect{7. Все ненулевые сдвиги и малые метки}
Качественное подавление распространяется и на сдвиги, стремящиеся к нулю на торе:
\[
                \sup_{v\ne0}\sup_{m\in\Z}
                   |N^{-1}\Tr(T_N(v)^* U_N^m)|=o(1).               \tag{12}
\]
Чтобы проверить равномерность, при каждом фиксированном $j\ge1$ используем
\[
 \limsup_{N\to\infty}\ \sup_{v\ne0}|(W_N^j)_{vv}|\le q^j,
                   \qquad q=e^{-0.8}.
\]
Действительно, при нарушении оценки можно выбрать последовательность $v_N$, а затем подпоследовательность, где либо $n(v_N)/N$ отделено от нуля, либо стремится к нулю. В первом случае применимо (2). Во втором случай фиксированного $j$ контролируется поднятой вейлевой эволюцией из раздела 6 части о жёсткости. В её взвешенной норме действует оценка $e^{-0.8 n(v_N)}q^{j-1}$ на диагональном элементе. Входной и выходной веса совпадают. Ошибка периодизации стремится к нулю при $n(v_N)=o(N)$, а всякая ненулевая метка имеет $n(v_N)\ge1$.

Рассмотрим спектральное среднее $A_0 T_N(v)$ из раздела 4 с любым фиксированным числом шагов $L$. В пределе сверху квадрат нормы не больше $(1+2q/(1-q))/L$, равномерно по $v\ne0$, по сумме геометрических диагональных корреляций. Это также верхняя граница для проекции на коммутант. Снимаем ограничение $L$, устремляя его к бесконечности после $N$. Поскольку все $U_N^m$ неподвижны для сопряжения и имеют норму 1, получаем (12).

\subt{Количественная малочастотная область} Для любого фиксированного $0<\epsilon<1$ имеем
\[
       \sup_{0<n(v)\le N^{1-\epsilon}}\sup_{m\in\Z}
                         |N^{-1}\Tr(T_N(v)^*U_N^m)|
                              =O_\epsilon((\log N)^{-1/2}).         \tag{13}
\]
Действительно, начнём с $R_0=\max(N^{1-\epsilon},\,100\log N)$, $R_i=8^iR_0$. При $0\le i\le L=\lfloor\epsilon\log N/(2\log 8)\rfloor$ все радиусы меньше $N/16$ для больших $N$. Из (3) части о жёсткости переход от $n(v)\le R_i$ за пределы $n(w)\le R_{i+1}$ имеет сумму модулей по поднятиям $O(e^{-3R_i})$ на столбец (при $n(Sv)<6R_i$ и расстоянии больше $2R_i$). Можно более грубо оценить норму всего хвоста через $O(N^2e^{-3R_i})$. Сравним по отдельности полную эволюцию и эволюцию, сжатую внутри $N/16$, с последовательными отсечениями по $R_i$. Унитарность или сжатие в каждом шаге даёт суммарную ошибку не больше $O(LN^2e^{-3R_0})$. На фиксированном ромбе периодизированная матрица отличается от поднятой в норме на $O(e^{-N})$ по (5) части о жёсткости. Таким образом, для всех $1\le j<L$ диагональный коэффициент ограничен $q^j+O(N^{-20})$ с одной константой, независимо от метки из указанной области. Спектральное среднее длины $L$ даёт $\|\Pi_{\{1\}}T_N(v)\|_{\rm n}^2=O(1/L)$ и (13).
\newpage
\sect{8. Частичные линейные возвраты порядка $N$}
Для любого целого $s$ обозначим $I_s=\operatorname{im}(S^s-I)\bmod N$. Разложение $M_N^s$ по модульным операторам Вейля имеет носитель ровно $I_s$, и все модули коэффициентов равны $|I_s|^{-1/2}$. Действительно, если $e\in\ker(S^s-I)$, то $T_N(e)$ коммутирует с $M_N^s$. Ненулевая метка его разложения должна лежать в симплектическом аннуляторе ядра. Он равен $I_s$: включение следует из симплектичности $S^s$, а обе мощности равны $N^2/|\ker(S^s-I)|$. Этот аргумент работает и на составном модуле. Точная ковариантность $M_N^sT_N(e)=T_N(S^se)M_N^s$ затем сдвигает модуль коэффициента на $(S^s-I)e$. Норма $\|M_N^s\|_{\rm n}=1$ даёт указанное значение.

Рассмотрим $N=7^k$ и фиксированное целое $a\ge1$. При $k>a$ выберем
\[
                   s=3u7^{k-a-1},\qquad u\ge1,\quad 7\nmid u .
\]
Тогда $I_s=7^{k-a}(\Z/N\Z)^2$ и $|I_s|=7^{2a}$. Для проверки $S^3-I=14\left(\begin{smallmatrix}7&10\\5&7\end{smallmatrix}\right)$; матрица внутри имеет определитель $-1$. Возведение $I+7^b C$ в степень 7 при $b\ge1$ даёт $I+7^{b+1}(C+7D)$ с целой матрицей $D$, по биному. Возведение в степень $u$, взаимно простую с 7, умножает $C\bmod7$ на $u$. Значит, $S^s-I=7^{k-a}C_s$, где $C_s$ обратима по модулю 7. Формула для носителя следует для всех таких $s$, включая растущие $u$.

У ненулевых меток $I_s$ величина $n(v)$ не меньше $N/7^a$. Поэтому для любого $m\in\Z$ из (8) получаем
\[
 \frac{|\Tr((M_N^s)^* U_N^m)|}{N}
     \le 7^{-a}\,\frac{|\Tr U_N^m|}{N}
                    +(1+o(1))\,T^{-1/2}
     \le 7^{-a}+O_a(T^{-1/2}).                                     \tag{14}
\]
Вклад нулевой метки выделен с точным модулем коэффициента $7^{-a}$; для остальных использованы Коши и (8) на целом пакете. В частности, при $u=1$ и $m=s=3N/7^{a+1}$ получаем строгий разрыв до частичного линейного возврата при времени порядка $N$:
\[
 \liminf_{k\to\infty}
         \min_{|\zeta|=1}\|U_{7^k}^{\,s}-\zeta M_{7^k}^{\,s}\|_{\rm n}
                                  \ge\sqrt{2(1-7^{-a})}.
\]
При $a=1$ это $\sqrt{12/7}$ на времени $3N/49$. Для $a=0$ чистая степень скалярна, её вейлев носитель равен $\{0\}$. Правая часть (14) тогда даёт лишь обычную тривиальную границу для $\Tr U_N^m$.

Если домножить линейную степень $M_N^s$ на $T_N(v_0)$ при $v_0\notin I_s$, её носитель лежит в смежном классе $v_0+I_s$, не содержащем нуля. При фиксированном $a$ размер этого класса постоянен по $N$; по (12) и Коши нормированный след $\Tr((T_N(v_0)M_N^s)^*U_N^m)/N$ стремится к нулю, равномерно по $m,s,v_0$ в этих диапазонах. Таким образом, также подавлены все сдвинутые частичные возвраты, чьи носители избегают нулевой метки.
\end{document}
